% SPDX-License-Identifier: PMPL-1.0-or-later
% Copyright (c) 2026 Jonathan D.A. Jewell (hyperpolymath) <j.d.a.jewell@open.ac.uk>
%
% typed-wasm: Applying Database Query Type Safety to WebAssembly Linear Memory
% Prepared for arXiv submission.

\documentclass[11pt,a4paper]{article}

%%% Packages %%%
\usepackage[utf8]{inputenc}
\usepackage[T1]{fontenc}
\usepackage{lmodern}
\usepackage{amsmath,amssymb,amsthm}
\usepackage{mathtools}
\usepackage{mathpartir}
\usepackage{listings}
\usepackage{booktabs}
\usepackage{tabularx}
\usepackage{hyperref}
\usepackage[margin=2.5cm]{geometry}
\usepackage{xcolor}
\usepackage{enumitem}
\usepackage{microtype}

%%% Theorem-like environments %%%
\newtheorem{definition}{Definition}
\newtheorem{theorem}{Theorem}
\newtheorem{conjecture}{Conjecture}

%%% Listings configuration %%%
\definecolor{codegreen}{rgb}{0,0.5,0}
\definecolor{codegray}{rgb}{0.5,0.5,0.5}
\definecolor{codepurple}{rgb}{0.58,0,0.82}
\definecolor{backcolour}{rgb}{0.97,0.97,0.97}

\lstdefinelanguage{Idris}{
  keywords={data,where,Type,let,in,do,if,then,else,total,covering,partial,
            import,module,public,export,private,mutual,namespace,record,
            interface,implementation,using,with,proof,impossible,case,of},
  sensitive=true,
  morecomment=[l]{--},
  morecomment=[s]{\{-}{-\}},
  morestring=[b]",
}

\lstset{
  backgroundcolor=\color{backcolour},
  basicstyle=\ttfamily\small,
  keywordstyle=\color{codepurple}\bfseries,
  commentstyle=\color{codegreen},
  stringstyle=\color{codegray},
  breaklines=true,
  frame=single,
  framesep=4pt,
  xleftmargin=6pt,
  xrightmargin=6pt,
  aboveskip=10pt,
  belowskip=10pt,
  numbers=none,
  captionpos=b,
}

%%% Custom macros %%%
\newcommand{\typedwasm}{\textsc{typed-wasm}}
\newcommand{\wasm}{Wasm}
\newcommand{\ie}{i.e.}
\newcommand{\eg}{e.g.}

%%% Inference rule helper (horizontal line with label) %%%
\newcommand{\levelrule}[1]{\textsc{#1}}

%%% Title %%%
\title{\typedwasm{}: Applying Database Query Type Safety\\to WebAssembly Linear Memory}
\author{
  Jonathan D.A.\ Jewell\\
  The Open University, Milton Keynes, UK\\
  \texttt{j.d.a.jewell@open.ac.uk}
}
\date{March 2026 --- Draft}

\begin{document}

\maketitle

\begin{abstract}
WebAssembly (Wasm) linear memory is an untyped byte array shared across
module boundaries.  When independently compiled modules --- potentially from
different source languages --- read and write the same memory regions, no
existing type system covers the cross-module interface.  We present
\typedwasm{}, a type system that applies a 10-level type safety
framework, originally developed for database query languages, to Wasm
linear memory.  The system treats contiguous memory segments as \emph{typed
region schemas} and load/store operations as \emph{typed projections}
verified against those schemas at compile time.  We formalise the system
in Idris~2 using Quantitative Type Theory (QTT), providing proofs of
bounds safety, aliasing freedom, effect purity, lifetime validity, and
linearity --- all erased before code generation, yielding zero runtime
overhead.  Our principal contribution is \textbf{multi-module schema
agreement}: a static verification that independently compiled Wasm
modules agree on the layout, types, alignment, and invariants of shared
memory regions --- a property that no source-level type system, and no
existing Wasm proposal, can express.
\end{abstract}

\medskip
\noindent\textbf{Keywords:}
WebAssembly, type safety, dependent types, linear memory,
multi-module verification, quantitative type theory, memory safety

\medskip
\noindent\textbf{ACM CCS:}
D.3.1 [Programming Languages]: Formal Definitions and Theory;
D.2.4 [Software Engineering]: Software/Program Verification;
F.3.1 [Logics and Meanings of Programs]: Specifying and Verifying Programs

\newpage
\tableofcontents
\newpage

%% ===================================================================
\section{Introduction}
\label{sec:introduction}

The WebAssembly specification~\cite{haas2017} defines a stack-based
virtual machine with strong instruction-level type safety: the validator
rejects programs whose stack operations are type-incorrect, and the
structured control flow ensures that branches target valid labels.  These
properties make Wasm a safe compilation target for individual languages.

However, Wasm's linear memory --- the flat byte array that serves as the
heap for compiled C, Rust, and other systems languages --- is entirely
untyped.  An \lstinline{i32.load offset=0} and an
\lstinline{f64.load offset=0} targeting the same address are both valid
instructions, regardless of what was actually stored there.  The Wasm
validator checks that the instruction encoding is correct, not that the
memory access is semantically meaningful.

This gap becomes a safety hazard when multiple independently compiled
modules share linear memory --- an increasingly common pattern in plugin
architectures, game engines, and polyglot Wasm compositions.  If Module~A
(compiled from Rust) writes a struct to shared memory and Module~B
(compiled from C or ReScript) reads the same memory with a different
layout assumption, the result is silent data corruption undetectable by
either source-level type system.

We observe that this problem has the same structure as untyped database
access: programs issue ``queries'' (load/store instructions) against a
``schema'' (memory layout) with no static guarantee that the query is
compatible with the schema.  The database query language community has
developed mature solutions to this problem, culminating in progressive
type safety frameworks that verify queries against schemas at compile time.

This paper presents \typedwasm{}, which applies this approach to Wasm
linear memory.  Our contributions are:

\begin{enumerate}
  \item \textbf{The database--memory analogy as a design principle}
    (Section~\ref{sec:analogy}): a structural correspondence between
    database schemas and memory region declarations, between queries and
    typed memory access, that yields a concrete grammar and type system.

  \item \textbf{A 10-level progressive type safety framework for Wasm
    memory} (Section~\ref{sec:levels}): adapting established levels
    (parse-time validity, schema binding, type compatibility, null
    safety, bounds proofs, result typing) and research levels (aliasing
    safety, effect tracking, lifetime safety, linearity) from the
    database domain to the memory domain.

  \item \textbf{Multi-module schema agreement}
    (Section~\ref{sec:multimodule}): a static verification that
    independently compiled Wasm modules agree on shared memory layout ---
    the principal contribution, addressing a class of bugs that no
    existing tool can detect.

  \item \textbf{A formalisation in Idris~2 with QTT}
    (Section~\ref{sec:formal}): dependent types that prove safety
    properties at compile time, with proof erasure yielding identical
    runtime code to hand-written Wasm.
\end{enumerate}

%% ===================================================================
\section{Problem Statement}
\label{sec:problem}

\subsection{The Untyped Byte Array}
\label{sec:untyped}

Wasm's linear memory model~\cite{wasmspec2024} exposes a resizable byte
array to all code within an instance.  Memory access instructions
(\lstinline{i32.load}, \lstinline{f64.store}, etc.)\ take an integer
offset and an alignment hint.  The Wasm validator ensures that the
alignment is valid for the instruction and that the access falls within
the allocated page range, but it imposes no constraints on the
\emph{semantic content} at the accessed offset.

Formally, if $M$ is the linear memory (a function from byte offsets to
bytes), then \lstinline{i32.load offset=k} computes
\[
  M[k] \;\big|\; \bigl(M[k{+}1] \ll 8\bigr) \;\big|\;
  \bigl(M[k{+}2] \ll 16\bigr) \;\big|\;
  \bigl(M[k{+}3] \ll 24\bigr)
\]
regardless of what was stored at offsets $k$ through $k{+}3$.  The
instruction cannot fail for type reasons --- only for out-of-bounds
access beyond the memory's page limit.

\subsection{The Cross-Module Boundary Problem}
\label{sec:crossmodule}

When multiple modules share linear memory, each module compiles with its
own struct layout assumptions.  Consider:

\begin{itemize}
  \item \textbf{Module~A} (Rust, \lstinline{#[repr(C)]}):
    \lstinline{Player \{ hp: i32 @ 0, _pad: [u8;4] @ 4, speed: f64 @ 8, name: [u8;32] @ 16 \}}
    --- 48~bytes, 8-byte aligned.

  \item \textbf{Module~B} (C, packed):
    \lstinline{player_t \{ int hp @ 0, double speed @ 4, char name[32] @ 12 \}}
    --- 44~bytes, unaligned.
\end{itemize}

Module~B reads \lstinline{speed} from offset~4, which in Module~A's
layout is padding.  Module~B writes \lstinline{name} starting at
offset~12, which overlaps Module~A's \lstinline{speed} field.  Neither
module's type checker detects this: Rust's borrow checker validates
within Rust; C's type system validates within C.  The mismatch exists at
the Wasm level, below both.

\subsection{Bug Taxonomy}
\label{sec:bugs}

We identify five classes of cross-module memory bugs:

\begin{table}[htbp]
\centering
\caption{Cross-module memory bug taxonomy.}
\label{tab:bugs}
\begin{tabularx}{\textwidth}{@{}lXX@{}}
\toprule
\textbf{Class} & \textbf{Description} & \textbf{Consequence} \\
\midrule
B1. Type reinterpretation & Load type differs from store type at same offset & Silent value corruption \\
B2. Layout disagreement & Modules assume different field offsets & All fields after the first are misread \\
B3. Null sentinel divergence & Modules disagree on the null representation & Valid pointers treated as null, or vice versa \\
B4. Lifetime mismatch & One module frees memory another still references & Use-after-free, data from reallocated memory \\
B5. Ownership confusion & Multiple modules believe they own the allocation & Double-free, heap corruption \\
\bottomrule
\end{tabularx}
\end{table}

These bugs share a common root: the absence of a shared schema that all
modules agree upon and that a static checker can verify.

\subsection{Insufficiency of Existing Approaches}
\label{sec:existing}

\begin{table}[htbp]
\centering
\caption{Existing approaches and their coverage gaps.}
\label{tab:existing}
\begin{tabularx}{\textwidth}{@{}lXX@{}}
\toprule
\textbf{Approach} & \textbf{What it covers} & \textbf{What it misses} \\
\midrule
Rust borrow checker & Ownership within Rust & Cross-module access from non-Rust code \\
Wasm validation & Instruction encoding, page bounds & Semantic correctness within pages \\
Wasm GC proposal & Managed reference types & Linear memory (unmanaged) \\
TAL~\cite{morrisett1999} & Register/stack types & Structured regions, cross-module schemas \\
MSWasm~\cite{disselkoen2022} & Segment bounds checking & Field-level types, schema agreement \\
Component Model~\cite{wasmcomponent2024} & Function call interfaces & Shared mutable memory \\
\bottomrule
\end{tabularx}
\end{table}

%% ===================================================================
\section{The Database--Memory Analogy}
\label{sec:analogy}

\subsection{Structural Correspondence}
\label{sec:correspondence}

The TypedQL framework~\cite{jewell2026} provides 10 progressive levels
of type safety for database query languages.  Its central mechanism is
verifying \emph{query projections} against a \emph{declared schema} at
compile time.  We observe that Wasm memory access has identical
structure:

\begin{table}[htbp]
\centering
\caption{Structural correspondence between databases and Wasm memory.}
\label{tab:correspondence}
\begin{tabularx}{\textwidth}{@{}lll@{}}
\toprule
\textbf{Database} & \textbf{Wasm Memory} & \textbf{Structural Role} \\
\midrule
Table & Region & Named container of typed records \\
Row & Region instance & Single record at a computed offset \\
Column & Field & Named, typed datum within a record \\
Schema & Region declaration & Compile-time layout specification \\
\lstinline{SELECT col FROM table} & \lstinline{region.get \$r .field} & Read projection \\
\lstinline{UPDATE table SET col = v} & \lstinline{region.set \$r .field, v} & Write projection \\
Foreign key & \lstinline{ptr<@OtherRegion>} & Cross-region typed reference \\
\lstinline{CHECK} constraint & \lstinline{invariant \{ ... \}} & Domain validity rule \\
\bottomrule
\end{tabularx}
\end{table}

\subsection{Why the Analogy Is Load-Bearing}
\label{sec:loadbearing}

This is not a surface-level metaphor.  The correspondence determines
concrete design decisions:

\begin{enumerate}
  \item \textbf{Region declarations} use declarative syntax isomorphic to
    \lstinline{CREATE TABLE}:
    \lstinline{region Players[100] \{ hp: i32; speed: f64; align 8; \}}.

  \item \textbf{Import/export of region schemas} mirrors shared database
    schemas across services.  Module~A exports; Module~B imports; the
    checker verifies structural compatibility.

  \item \textbf{Cross-region invariants} are foreign key constraints over
    memory:
    \lstinline{invariant \{ across: Enemies, Players; holds: forall e. 0 <= e.target\_id < 4096; \}}.

  \item \textbf{\lstinline{region.scan} with \lstinline{where}} is
    \lstinline{SELECT * FROM region WHERE predicate}, extending naturally
    to iteration, filtering, and aggregation over array regions.
\end{enumerate}

The 10-level framework transfers directly because the safety properties
are properties of \emph{projections against schemas}, independent of
whether the schema describes a database table or a memory region.

\subsection{Where the Analogy Diverges}
\label{sec:divergences}

Three aspects of Wasm memory have no direct database counterpart:

\begin{itemize}
  \item \textbf{Alignment and padding.}  Memory fields require alignment
    to natural boundaries.  \typedwasm{} makes alignment explicit in
    declarations and includes it in the type-checking contract.

  \item \textbf{Pointer arithmetic.}  Wasm linear memory is addressed by
    integer offsets.  \typedwasm{} introduces pointer types
    (\lstinline{ptr}, \lstinline{ref}, \lstinline{unique}) with ownership
    semantics borrowed from substructural type theory.

  \item \textbf{Concurrency.}  Database transactions provide isolation;
    Wasm shared memory provides none.  \typedwasm{}'s aliasing safety
    (Level~7) addresses single-threaded aliasing; concurrent access is
    deferred to future work.
\end{itemize}

%% ===================================================================
\section{The 10 Levels of Type Safety for Wasm Memory}
\label{sec:levels}

TypeLL defines a progressive stack of 10 type safety levels.  Each level
addresses a specific class of failure.  Levels are cumulative: a program
cannot achieve Level~$n$ without satisfying Levels~1 through $n{-}1$.
Simple operations exit early --- a read from a singleton region achieves
Level~6 with no ownership, effect, lifetime, or linearity analysis.

\subsection{Established Levels (1--6)}
\label{sec:established}

\paragraph{Level 1 --- Instruction Validity.}
The access expression conforms to the \typedwasm{} grammar.  Analogous to
parse-time safety for queries.  Wasm itself validates instruction
encoding but not higher-level access structure.

\paragraph{Level 2 --- Region-Binding.}
Every \lstinline{region.get \$target .field} resolves \lstinline{\$target}
to a declared region and \lstinline{.field} to a declared field.
Analogous to schema binding.  Wasm has no concept of named regions or
fields.

\paragraph{Level 3 --- Type-Compatible Access.}
The result type of a memory access is determined by the field
declaration, not by the programmer's choice of load instruction.  If
\lstinline{hp: i32}, then \lstinline{region.get \$p .hp} produces
\lstinline{i32} and compiles to \lstinline{i32.load} --- never
\lstinline{f64.load}.  Analogous to type-compatible operations in
queries.

\paragraph{Level 4 --- Null Safety.}
The type system distinguishes \lstinline{ptr<T>} (non-nullable) from
\lstinline{opt<ptr<T>>} (nullable).  Accessing a nullable value requires
explicit unwrapping.  Analogous to null safety in query result handling.

\paragraph{Level 5 --- Bounds-Proof.}
Every memory access is provably within the region's bounds.  For static
indices, the Idris~2 prover verifies the bound at compile time.  For
dynamic indices, a minimal runtime check is inserted only where the
prover cannot discharge the obligation.  Analogous to injection-proof
safety (separation of structure from data).

\paragraph{Level 6 --- Result-Type.}
The return type of every memory access expression is statically known and
propagated through bindings and function returns.  Analogous to
result-type safety in queries.

\subsection{Research Levels (7--10)}
\label{sec:research}

\paragraph{Level 7 --- Aliasing Safety.}
Mutable references (\lstinline{&mut}) to the same region do not alias.
\typedwasm{} provides three pointer kinds:

\begin{itemize}
  \item \lstinline{&} (shared borrow): unlimited concurrent readers, no
    mutation.
  \item \lstinline{&mut} (exclusive borrow): exactly one mutable
    reference, no concurrent readers.
  \item \lstinline{own} (owning): single owner, must transfer or free
    (Level~10).
\end{itemize}

The prover verifies exclusivity at every program point.  This is
structurally analogous to Rust's borrow discipline, but operates at the
Wasm region schema level rather than the source language level.

\paragraph{Level 8 --- Effect Tracking.}
Functions declare their memory effects:
\lstinline{effects \{ ReadRegion(Players), WriteRegion(Config) \}}.
A function declared read-only cannot perform writes.  Effects are checked
by subsumption: $\mathit{actual} \subseteq \mathit{declared}$.  Analogous
to distinguishing \lstinline{SELECT} from \lstinline{UPDATE} in queries.

\paragraph{Level 9 --- Lifetime Safety.}
References carry lifetime annotations scoping their validity.  A
reference with lifetime $\ell$ is valid only while $\ell$ is live.  The
prover verifies that no reference is used after the memory it points to
has been freed.  Analogous to temporal freshness in queries.

\paragraph{Level 10 --- Linearity.}
Owning handles are bound with QTT quantity~1, enforcing exactly-once
consumption.  An allocation must be freed exactly once: the type checker
rejects programs with zero uses (leak) or two uses (double-free) of the
handle.  This composes with Level~9: freeing invalidates the lifetime,
making all references with that lifetime unusable.  Analogous to
connection linearity in database access.

\subsection{Cross-Level Composition}
\label{sec:composition}

The levels compose to provide compound guarantees:

\begin{itemize}
  \item \textbf{L5 $\wedge$ L7:} Bounds-proof plus aliasing safety
    proves that two exclusive references never overlap in memory, even
    for different fields within the same region.

  \item \textbf{L8 $\wedge$ L9:} Effect tracking plus lifetime safety
    proves that a read-only reference cannot observe a free effect on its
    target region.

  \item \textbf{L7 $\wedge$ L10:} Aliasing safety plus linearity yields
    the full ownership model at the Wasm memory level.

  \item \textbf{L2 $\wedge$ L3 $\wedge$ multi-module:} Region-binding
    plus type-compatibility across module boundaries is the composition
    unique to \typedwasm{}.
\end{itemize}

%% ===================================================================
\section{Multi-Module Schema Agreement}
\label{sec:multimodule}

\subsection{The Verification Problem}
\label{sec:verification}

Given Module~A exporting region $R_A$ with schema $S_A$ and Module~B
importing region $R_B$ (claiming to correspond to $R_A$) with schema
$S_B$, verify that $S_B$ is structurally compatible with $S_A$.

\subsection{Compatibility Relation}
\label{sec:compatibility}

We define structural compatibility $S_B \preceq S_A$ (\emph{import
compatible with export}) as follows:

\begin{definition}[Field Compatibility]
\label{def:fieldcompat}
Fields $f_A$ and $f_B$ are compatible, written $f_B \sim f_A$, iff:
\begin{itemize}
  \item $f_B.\mathit{name} = f_A.\mathit{name}$
  \item $f_B.\mathit{type} = f_A.\mathit{type}$ (strict equality, no
    implicit conversions)
  \item $f_B.\mathit{offset} = f_A.\mathit{offset}$ (computed from types
    and alignment)
\end{itemize}
\end{definition}

\begin{definition}[Schema Compatibility]
\label{def:schemacompat}
$S_B \preceq S_A$ iff:
\begin{itemize}
  \item For every field $f_B \in S_B$, there exists $f_A \in S_A$ such
    that $f_B \sim f_A$.
  \item $S_B.\mathit{alignment} = S_A.\mathit{alignment}$
  \item $S_B.\mathit{instance\_count} = S_A.\mathit{instance\_count}$
\end{itemize}

Note that $\preceq$ permits $S_B$ to be a \emph{subset} of $S_A$:
Module~B may import fewer fields than Module~A exports.  Module~B simply
cannot access the omitted fields.  This is analogous to a database view
that projects a subset of columns.
\end{definition}

\begin{definition}[Multi-Module Agreement]
\label{def:multimodule}
A set of modules $\{M_1, \ldots, M_n\}$ sharing region $R$ satisfies
multi-module agreement iff exactly one module $M_e$ exports $R$ with
schema $S_e$, and for every importing module $M_i$ with schema $S_i$,
$S_i \preceq S_e$.
\end{definition}

\subsection{Formalisation}
\label{sec:formalisation}

In the Idris~2 ABI layer, schema agreement is a dependent type:

\begin{lstlisting}[language=Idris]
data CompatCertificate : Type where
  MkCompat : (agreement : SchemaAgreement)
          -> (alignment : AlignmentAgrees ea ia)
          -> (instances : InstanceCountAgrees ec ic)
          -> CompatCertificate
\end{lstlisting}

Constructing a \lstinline{CompatCertificate} requires providing proofs
of each component.  If any proof cannot be constructed --- a field name
mismatch, a type difference, an alignment disagreement --- the
certificate cannot be produced, and the program is rejected at compile
time.

\subsection{Diagnostic Reporting}
\label{sec:diagnostics}

When verification fails, the checker produces structured diagnostics:

\begin{lstlisting}[basicstyle=\ttfamily\small,language={}]
Schema mismatch for region 'Player' imported by module_b from module_a:
  Field 'speed': type mismatch
    Expected (from module_a): f64
    Found (in module_b):      f32
  Field 'name': offset mismatch
    Expected (from module_a): offset 16 (8-byte aligned)
    Found (in module_b):      offset 12 (assumes packed layout)
\end{lstlisting}

\subsection{Practical Applications}
\label{sec:applications}

Multi-module schema agreement enables:

\begin{itemize}
  \item \textbf{Plugin architectures} where the host exports region
    schemas and plugins import them.  Schema-breaking changes are caught
    at plugin compile time.

  \item \textbf{Polyglot Wasm compositions} where Rust, C, and ReScript
    modules share structured data with compile-time layout verification.

  \item \textbf{Hot-reloadable modules} where a new module version is
    checked for schema compatibility with existing modules before
    deployment.
\end{itemize}

%% ===================================================================
\section{Implications for GraalVM}
\label{sec:graalvm}

GraalVM's Truffle framework~\cite{wurthinger2017} enables cross-language
interoperation through \lstinline{InteropLibrary}, which dispatches
member reads, writes, and invocations at runtime.  When a JavaScript
object is accessed from Python code, \lstinline{InteropLibrary.readMember()}
performs runtime type checking and conversion.

This has two costs: runtime dispatch overhead and the absence of
compile-time guarantees.  \typedwasm{}'s region schemas could serve as a
compile-time contract between Truffle languages:

\begin{enumerate}
  \item Language~A exports a region schema describing shared state.
  \item Language~B imports the schema.
  \item The checker verifies agreement at AOT compilation time.
  \item The Truffle partial evaluator replaces \lstinline{InteropLibrary}
    dispatch with direct memory access at known offsets, justified by the
    schema proof.
\end{enumerate}

This applies specifically to GraalVM's native interop paths (Sulong,
shared memory).  Adapting the approach to Truffle's object-based interop
for managed languages requires mapping object shapes to region schemas,
which is feasible given Truffle's \lstinline{DynamicObject} shape
tracking but remains future work.

%% ===================================================================
\section{Formal Foundations}
\label{sec:formal}

\subsection{Idris 2 and Dependent Types}
\label{sec:idris}

\typedwasm{}'s proofs are written in Idris~2~\cite{brady2021}, a
dependently typed language where types may depend on values.  This is
necessary because memory safety properties inherently depend on values:
``index~$i$ is within the bounds of array of size~$n$'' is a proposition
over the value of~$i$.

A region schema is a value-level record:

\begin{lstlisting}[language=Idris]
data Schema : List (String, FieldType) -> Type where
  Nil  : Schema []
  (::) : Field name ty -> Schema rest -> Schema ((name, ty) :: rest)
\end{lstlisting}

A typed access is indexed by the schema:

\begin{lstlisting}[language=Idris]
data TypedGet : Schema fields -> (name : String) -> FieldType -> Type where
  GetHere  : TypedGet (MkField {name} {ty} :: rest) name ty
  GetThere : TypedGet rest name ty -> TypedGet (f :: rest) name ty
\end{lstlisting}

The result type of \lstinline{region.get \$r .name} is \emph{computed}
from the schema, not declared by the programmer.

\subsection{Quantitative Type Theory}
\label{sec:qtt}

Idris~2's QTT~\cite{atkey2018,mcbride2016} assigns quantities to
bindings:

\begin{table}[htbp]
\centering
\caption{QTT quantities and their application in \typedwasm{}.}
\label{tab:qtt}
\begin{tabularx}{\textwidth}{@{}llX@{}}
\toprule
\textbf{Quantity} & \textbf{Usage} & \textbf{\typedwasm{} application} \\
\midrule
$0$ (erased) & Compile-time only & Proof witnesses, schema metadata, bounds evidence \\
$1$ (linear) & Exactly once at runtime & Owning handles, allocation tokens \\
$\omega$ (unrestricted) & Any number of times & Normal values, shared references \\
\bottomrule
\end{tabularx}
\end{table}

Level~10 is implemented by binding allocation handles with quantity~1:

\begin{lstlisting}[language=Idris]
freeRegion : (1 _ : LinHandle token) -> RegionLive token -> FreeResult
\end{lstlisting}

The \lstinline{(1 _)} annotation means the type checker rejects any
program where the handle is used zero times (leak) or more than once
(double-free).

\subsection{Proof Erasure}
\label{sec:erasure}

All quantity-0 terms --- which includes all proof witnesses, bounds
evidence, schema metadata, compatibility certificates, and lifetime
annotations --- are erased by the Idris~2 compiler before code
generation.

\begin{theorem}[Proof Erasure, informal]
\label{thm:erasure}
The runtime representation of a well-typed \typedwasm{} program is
identical to the Wasm instructions that a hand-written program would
produce.  The type safety is a compile-time property with no runtime
manifestation.
\end{theorem}

This is the zero-overhead guarantee: \lstinline{region.get \$player .hp}
compiles to exactly \lstinline{i32.load offset=0}, with the proof that
this instruction is correct existing only during compilation.

\subsection{Soundness Sketch}
\label{sec:soundness}

We conjecture the following soundness property:

\begin{conjecture}[\typedwasm{} Soundness]
\label{conj:soundness}
If a \typedwasm{} program type-checks at all 10 levels, then the
compiled Wasm program:
\begin{enumerate}[label=(\alph*)]
  \item never performs a type-confused memory read (L3),
  \item never accesses memory outside a declared region (L5),
  \item never holds aliased mutable references (L7),
  \item never performs an undeclared memory effect (L8),
  \item never dereferences a freed region (L9), and
  \item never leaks or double-frees an allocation (L10).
\end{enumerate}
\end{conjecture}

A full mechanised proof is future work.  The current Idris~2
formalisation provides dependent-type-level evidence for
properties~(a)--(f), verified by the totality checker, but does not yet
connect to a formal Wasm operational semantics.

%% ===================================================================
\section{Implementation Architecture}
\label{sec:implementation}

\begin{table}[htbp]
\centering
\caption{Implementation layers of \typedwasm{}.}
\label{tab:implementation}
\begin{tabularx}{\textwidth}{@{}llX@{}}
\toprule
\textbf{Layer} & \textbf{Language} & \textbf{Role} \\
\midrule
ABI definitions & Idris~2 & Dependent types: region schemas, access rules, proof obligations \\
FFI bridge & Zig & C-ABI runtime: region management, typed load/store, Wasm codegen \\
Surface parser & ReScript & Parse \texttt{.twasm} syntax to typed AST \\
Proof engine & Idris~2 totality checker & Discharge Levels~5--10 obligations \\
Code generator & Zig $\to$ Wasm & Emit raw Wasm instructions from verified typed access \\
Ecosystem plugin & Rust (TypedQLiser) & Integration as a TypedQLiser target language \\
\bottomrule
\end{tabularx}
\end{table}

\paragraph{Zig FFI.}
Zig provides native C ABI compatibility, built-in
\lstinline{wasm32-freestanding} target support, cross-compilation without
runtime dependencies, and memory-safe defaults.  Region handles encode
base offset, schema ID, generation counter (for lifetime tracking), and
ownership flag in a 64-bit value.

\paragraph{ReScript parser.}
The surface syntax parser is written in ReScript, consistent with the
VQL-UT parser in the TypedQL ecosystem.  ReScript compiles to
JavaScript/Wasm via Deno.

\paragraph{TypedQLiser integration.}
\typedwasm{} implements the \lstinline{QueryLanguagePlugin} trait from
TypedQLiser, making Wasm memory a target alongside SQL, GraphQL, and VQL.
This enables end-to-end verification: a database query is type-checked by
TypedQLiser, and the memory region holding its results is type-checked by
\typedwasm{}.

%% ===================================================================
\section{Related Work}
\label{sec:related}

\paragraph{Typed Assembly Language.}
Morrisett et al.~\cite{morrisett1999} introduced TAL, assigning types to
registers and stack slots to prove progress and preservation for assembly
programs.  \typedwasm{} extends this approach with structured region
schemas, cross-module import/export, and a 10-level progressive
framework.  TAL types individual memory cells; \typedwasm{} types
structured multi-field regions.

\paragraph{MSWasm.}
Disselkoen et al.~\cite{disselkoen2022} add memory segments with bounds
checking to Wasm, enforcing spatial safety.  \typedwasm{} builds on this
by adding field-level types within bounded regions, cross-module schema
agreement, effect tracking, and linearity.  The two are complementary:
MSWasm ensures access is within bounds; \typedwasm{} ensures access is
also type-correct within those bounds.

\paragraph{Wasm GC Proposal.}
The GC proposal~\cite{wasmgc2024} introduces managed struct and array
types.  These are type-safe by construction but apply only to GC-managed
objects, not to linear memory.  Systems languages (C, Rust) that manage
their own memory cannot benefit.  \typedwasm{} and Wasm GC are
complementary.

\paragraph{Wasm Component Model.}
The Component Model~\cite{wasmcomponent2024} types function call
interfaces with automatic marshalling.  It addresses isolation-based
composition (copying data across boundaries).  \typedwasm{} addresses
shared-memory composition (typing in-place access).  Both are valid
architectural choices; \typedwasm{} covers the case where copying is too
expensive.

\paragraph{Rust Ownership.}
Rust's affine type system~\cite{jung2018} prevents use-after-free,
double-free, and data races within Rust code.  \typedwasm{}'s Levels~7,
9, and~10 are deliberately analogous but operate at the Wasm level,
covering cross-language boundaries that Rust cannot reach.

\paragraph{Checked C and Typed LLVM.}
Checked~C~\cite{elliott2018} adds bounds-checked pointers to C.
Typed LLVM adds types to LLVM IR.  Both improve single-language
compilation pipelines.  \typedwasm{} operates at the Wasm convergence
point, where multiple source languages meet.

\paragraph{Substructural Type Systems.}
Linear types~\cite{wadler1990}, uniqueness types~\cite{barendsen1996},
and quantitative type theory~\cite{atkey2018} provide the theoretical
foundation for \typedwasm{}'s Levels~7--10.  Our contribution is not new
type theory but rather the application of these ideas to the specific
problem of cross-module Wasm memory safety.

%% ===================================================================
\section{Limitations and Future Work}
\label{sec:limitations}

\paragraph{Concurrency.}
\typedwasm{} v0.1 does not address Wasm shared memory with atomics.
Concurrent access to shared regions requires session types or fractional
permissions, which are explored in the TypeQL-Experimental project but
not yet integrated.

\paragraph{GC interaction.}
The boundary between linear memory and GC-managed objects is currently
untyped.  A \lstinline{gcref<T>} field type is conceivable but requires
runtime cooperation beyond static checking.

\paragraph{Dynamic layouts.}
Programs that determine memory layout at runtime (JIT compilers, dynamic
object models) cannot use static region declarations without an
adaptation layer.

\paragraph{Adoption barrier.}
All participating modules must carry \typedwasm{} annotations.  Automatic
schema inference from \lstinline{#[repr(C)]} attributes, C headers, and
Wasm debug custom sections could lower this barrier.

\paragraph{Proof completeness.}
Levels~5--10 depend on the Idris~2 totality checker.  Complex arithmetic
bounds may require manual proof assistance.  The practical frequency of
this in real Wasm modules is an empirical question.

\paragraph{Mechanised soundness.}
The soundness conjecture (Section~\ref{sec:soundness}) is not yet
mechanically verified against a formal Wasm operational semantics.
Connecting the Idris~2 formalisation to an existing Wasm mechanisation
(\eg, WasmCert~\cite{watt2018}) is a priority.

%% ===================================================================
\section{Conclusion}
\label{sec:conclusion}

Wasm's linear memory is the largest untyped shared mutable state surface
in modern software.  Multiple languages compile to Wasm; multiple modules
share linear memory; no existing type system covers the boundary between
them.

\typedwasm{} addresses this by recognising that Wasm memory access has
the same structure as database queries --- projections against a declared
schema.  By applying a 10-level type safety framework to memory regions,
\typedwasm{} provides:

\begin{enumerate}
  \item \textbf{Region schemas} that declare memory layout with field
    names, types, alignment, and invariants.

  \item \textbf{Multi-module schema agreement} that statically verifies
    layout compatibility across independently compiled modules.

  \item \textbf{10 levels of progressive type safety} from instruction
    validity through linearity, with formal proofs in Idris~2 and zero
    runtime overhead through proof erasure.

  \item \textbf{A practical architecture} integrating Idris~2, Zig,
    ReScript, and the TypedQLiser ecosystem.
\end{enumerate}

The implication extends beyond Wasm.  Any system where multiple languages
share untyped mutable state --- GraalVM Truffle interop, JNI, FFI
bridges, shared memory IPC --- faces the same class of bugs.
\typedwasm{} demonstrates that database query type safety is a
transferable framework for these domains.

%% ===================================================================
\appendix
\section{Formal Semantics: Inference Rules}
\label{app:inference}

This appendix presents the formal typing rules for each of the 10 levels
of \typedwasm{}.  We write $\Gamma$ for the variable context, $\Sigma$
for the region store (mapping region names to schemas), and $S$ for a
schema (mapping field names to types and offsets).

\subsection{Level 1: Instruction Validity}

An access expression is valid if the region and field names are
syntactically well-formed:

\begin{mathpar}
\inferrule*[left=\levelrule{L1-Valid}]
  {r \in \operatorname{dom}(\Sigma) \\ f \in \operatorname{fields}(\Sigma(r))}
  {\Gamma \vdash \mathsf{region.get}\;r\;.f : \mathsf{valid}}
\end{mathpar}

\subsection{Level 2: Region Binding}

Every access resolves to a concrete region and field within the store:

\begin{mathpar}
\inferrule*[left=\levelrule{L2-Bind}]
  {\Sigma(r) = S \\ f \in S}
  {\Gamma;\, \Sigma \vdash \mathsf{region.get}\;r\;.f \leadsto (r, S, f)}
\end{mathpar}

\subsection{Level 3: Type Compatibility}

The result type of a memory access is determined by the schema, not by
the programmer's choice of instruction:

\begin{mathpar}
\inferrule*[left=\levelrule{L3-Type}]
  {\Gamma;\, \Sigma \vdash \mathsf{region.get}\;r\;.f \leadsto (r, S, f) \\
   S(f) = \tau}
  {\Gamma;\, \Sigma \vdash \mathsf{region.get}\;r\;.f : \tau}
\end{mathpar}

\subsection{Level 4: Null Safety}

Nullable pointer access requires explicit case analysis:

\begin{mathpar}
\inferrule*[left=\levelrule{L4-Null}]
  {\Gamma;\, \Sigma \vdash e : \mathsf{opt}\langle\mathsf{ptr}\langle R \rangle\rangle \\
   \Gamma,\, x : \mathsf{ptr}\langle R \rangle \vdash e_1 : \tau \\
   \Gamma \vdash e_2 : \tau}
  {\Gamma;\, \Sigma \vdash \mathsf{match}\;e\;\{
     \mathsf{Some}(x) \to e_1,\;
     \mathsf{None} \to e_2
   \} : \tau}
\end{mathpar}

\subsection{Level 5: Bounds Proof}

Every indexed access carries compile-time evidence of being within bounds:

\begin{mathpar}
\inferrule*[left=\levelrule{L5-Bounds}]
  {\Sigma(r) = S \\ S.\mathit{count} = n \\ 0 \leq i < n}
  {\Gamma;\, \Sigma \vdash r[i] : \mathsf{InBounds}}
\end{mathpar}

\subsection{Level 6: Result Type}

The host-language type of every access is statically determined from the
schema:

\begin{mathpar}
\inferrule*[left=\levelrule{L6-Result}]
  {\Gamma;\, \Sigma \vdash \mathsf{region.get}\;r\;.f : \tau \\
   \tau = \mathsf{HostType}(S(f).\mathit{type})}
  {\Gamma;\, \Sigma \vdash \mathsf{region.get}\;r\;.f : \mathsf{HostType}(S(f).\mathit{type})}
\end{mathpar}

\subsection{Level 7: Aliasing Safety}

Mutable references require exclusivity.  Shared and exclusive borrows
are incompatible:

\begin{mathpar}
\inferrule*[left=\levelrule{L7-Shared}]
  {\Gamma;\, \Sigma \vdash e : \mathsf{ptr}\langle R \rangle \\
   R \notin \operatorname{mutBorrows}(\Gamma)}
  {\Gamma;\, \Sigma \vdash \mathsf{borrow}_{\&}\;e : \&R}

\inferrule*[left=\levelrule{L7-Exclusive}]
  {\Gamma;\, \Sigma \vdash e : \mathsf{ptr}\langle R \rangle \\
   R \notin \operatorname{borrows}(\Gamma)}
  {\Gamma;\, \Sigma \vdash \mathsf{borrow}_{\&\mathsf{mut}}\;e : \&\mathsf{mut}\;R}
\end{mathpar}

\subsection{Level 8: Effect Tracking}

Function effects are checked by subsumption --- the actual effects must
be a subset of the declared effects:

\begin{mathpar}
\inferrule*[left=\levelrule{L8-Effect}]
  {\Gamma;\, \Sigma \vdash_{\mathit{eff}} \mathit{body} : \mathit{actual} \\
   \mathit{actual} \subseteq \mathit{declared}}
  {\Gamma;\, \Sigma \vdash \mathsf{fn}\;(\mathit{declared})\;\mathit{body} : \mathit{ok}}

\inferrule*[left=\levelrule{L8-Read}]
  {\Gamma;\, \Sigma \vdash \mathsf{region.get}\;r\;.f : \tau}
  {\Gamma;\, \Sigma \vdash_{\mathit{eff}} \mathsf{region.get}\;r\;.f :
     \{\mathsf{ReadRegion}(r)\}}

\inferrule*[left=\levelrule{L8-Write}]
  {\Gamma;\, \Sigma \vdash \mathsf{region.set}\;r\;.f\;v : \mathsf{unit}}
  {\Gamma;\, \Sigma \vdash_{\mathit{eff}} \mathsf{region.set}\;r\;.f\;v :
     \{\mathsf{WriteRegion}(r)\}}
\end{mathpar}

\subsection{Level 9: Lifetime Safety}

References carry lifetime annotations; a reference is valid only while
its lifetime is live.  The outlives relation $\ell_1 \sqsupseteq \ell_2$
means $\ell_1$ is valid at least as long as $\ell_2$:

\begin{mathpar}
\inferrule*[left=\levelrule{L9-Ref}]
  {\Gamma;\, \Sigma \vdash e : \mathsf{ptr}\langle R \rangle \\
   \ell \in \operatorname{liveLifetimes}(\Gamma)}
  {\Gamma;\, \Sigma \vdash \mathsf{ref}^{\ell}\;e : \&^{\ell} R}

\inferrule*[left=\levelrule{L9-Outlives}]
  {\Gamma \vdash x : \&^{\ell_1} R \\
   \ell_2 \sqsupseteq \ell_1}
  {\Gamma \vdash x : \&^{\ell_2} R}

\inferrule*[left=\levelrule{L9-UseAfterFree}]
  {\Gamma;\, \Sigma \vdash \mathsf{free}\;r : \mathsf{unit} \\
   \ell_r = \operatorname{lifetime}(r)}
  {\ell_r \notin \operatorname{liveLifetimes}(\Gamma')}
\end{mathpar}

\subsection{Level 10: Linearity}

Owning handles use QTT quantity~1.  An owning handle must be consumed
exactly once:

\begin{mathpar}
\inferrule*[left=\levelrule{L10-Alloc}]
  {\Gamma;\, \Sigma \vdash \mathsf{alloc}\;S\;n : \mathsf{result}}
  {\Gamma;\, \Sigma \vdash \mathsf{alloc}\;S\;n :
     (1\;h : \mathsf{LinHandle}\;\mathit{tok}) \times
     \mathsf{RegionLive}\;\mathit{tok}}

\inferrule*[left=\levelrule{L10-Free}]
  {1\;h : \mathsf{LinHandle}\;\mathit{tok} \in \Gamma \\
   \mathsf{RegionLive}\;\mathit{tok} \in \Gamma}
  {\Gamma;\, \Sigma \vdash \mathsf{free}\;h : \mathsf{FreeResult}}

\inferrule*[left=\levelrule{L10-NoLeak}]
  {1\;h : \mathsf{LinHandle}\;\mathit{tok} \in \Gamma \\
   h \notin \operatorname{consumed}(\mathit{body})}
  {\Gamma;\, \Sigma \vdash \mathit{body} : \bot
   \quad\text{(rejected: linear variable not consumed)}}
\end{mathpar}

\subsection{Multi-Module Schema Agreement}

The compatibility judgement for cross-module schema verification:

\begin{mathpar}
\inferrule*[left=\levelrule{Schema-Compat}]
  {\forall\, f \in S_{\mathit{import}}.\;
     \exists\, f' \in S_{\mathit{export}}.\;
       f.\mathit{name} = f'.\mathit{name} \;\wedge\;
       f.\mathit{type} = f'.\mathit{type} \;\wedge\;
       f.\mathit{offset} = f'.\mathit{offset} \\
   S_{\mathit{import}}.\mathit{align} = S_{\mathit{export}}.\mathit{align} \\
   S_{\mathit{import}}.\mathit{count} = S_{\mathit{export}}.\mathit{count}}
  {S_{\mathit{import}} \preceq S_{\mathit{export}}}
\end{mathpar}

\begin{mathpar}
\inferrule*[left=\levelrule{Multi-Agree}]
  {M_e \;\mathsf{exports}\; R \;\mathsf{with}\; S_e \\
   \forall\, i \in \{1, \ldots, n\} \setminus \{e\}.\;
     M_i \;\mathsf{imports}\; R \;\mathsf{with}\; S_i \;\wedge\;
     S_i \preceq S_e}
  {\{M_1, \ldots, M_n\} \vdash R : \mathsf{agreed}}
\end{mathpar}

%% ===================================================================
\begin{thebibliography}{20}

\bibitem{atkey2018}
R.~Atkey.
\newblock Syntax and semantics of quantitative type theory.
\newblock In \emph{ACM/IEEE Symposium on Logic in Computer Science (LICS)},
  2018.

\bibitem{barendsen1996}
E.~Barendsen and S.~Smetsers.
\newblock Uniqueness typing for functional languages with graph rewriting
  semantics.
\newblock \emph{Mathematical Structures in Computer Science}, 6(6):579--612,
  1996.

\bibitem{brady2021}
E.~Brady.
\newblock Idris 2: Quantitative type theory in practice.
\newblock In \emph{European Conference on Object-Oriented Programming (ECOOP)},
  2021.

\bibitem{disselkoen2022}
C.~Disselkoen, J.~Renner, C.~Watt, T.~Garber, S.~Cauligi, and D.~Stefan.
\newblock {MSWasm}: Soundly enforcing memory-safe execution of unsafe code.
\newblock In \emph{IEEE Symposium on Security and Privacy}, 2022.

\bibitem{elliott2018}
A.~S. Elliott, A.~Ruef, M.~Hicks, and D.~Tarditi.
\newblock {Checked C}: Making {C} safe by extension.
\newblock In \emph{IEEE Cybersecurity Development (SecDev)}, 2018.

\bibitem{haas2017}
A.~Haas, A.~Rossberg, D.~L. Schuff, et~al.
\newblock Bringing the web up to speed with {WebAssembly}.
\newblock In \emph{ACM SIGPLAN Conference on Programming Language Design and
  Implementation (PLDI)}, 2017.

\bibitem{jewell2026}
J.~D.~A. Jewell.
\newblock {TypedQL}: A 10-level type safety framework for database query
  languages.
\newblock Technical report, hyperpolymath, 2026.

\bibitem{jung2018}
R.~Jung, J.-H. Jourdan, R.~Krebbers, and D.~Dreyer.
\newblock {RustBelt}: Securing the foundations of the {Rust} programming
  language.
\newblock \emph{Proceedings of the ACM on Programming Languages (POPL)},
  2(POPL):66, 2018.

\bibitem{mcbride2016}
C.~McBride.
\newblock I got plenty o' nuttin'.
\newblock In \emph{A List of Successes That Can Change the World}, LNCS
  9600:207--233, 2016.

\bibitem{morrisett1999}
G.~Morrisett, D.~Walker, K.~Crary, and N.~Glew.
\newblock From {System F} to typed assembly language.
\newblock \emph{ACM Transactions on Programming Languages and Systems},
  21(3):527--568, 1999.

\bibitem{wadler1990}
P.~Wadler.
\newblock Linear types can change the world!
\newblock In \emph{Programming Concepts and Methods}, 1990.

\bibitem{watt2018}
C.~Watt.
\newblock Mechanising and verifying the {WebAssembly} specification.
\newblock In \emph{ACM SIGPLAN Conference on Certified Programs and Proofs
  (CPP)}, 2018.

\bibitem{wasmgc2024}
{WebAssembly Community Group}.
\newblock {GC} proposal, 2024.
\newblock \url{https://github.com/WebAssembly/gc}.

\bibitem{wasmcomponent2024}
{WebAssembly Community Group}.
\newblock Component model, 2024.
\newblock \url{https://github.com/WebAssembly/component-model}.

\bibitem{wasmspec2024}
{WebAssembly Community Group}.
\newblock {WebAssembly} specification, 2024.
\newblock \url{https://webassembly.github.io/spec/}.

\bibitem{wurthinger2017}
T.~W\"urthinger, C.~Wimmer, C.~Humer, et~al.
\newblock Practical partial evaluation for high-performance dynamic language
  runtimes.
\newblock In \emph{ACM SIGPLAN Conference on Programming Language Design and
  Implementation (PLDI)}, 2017.

\end{thebibliography}

\bigskip
\noindent\rule{\textwidth}{0.4pt}

\smallskip
\noindent\emph{\typedwasm{} is part of the hyperpolymath TypeLL ecosystem.}\\
\noindent\emph{Repository: \url{https://github.com/hyperpolymath/typed-wasm}}

\end{document}
